\IfFileExists{stacks-project.cls}{%
\documentclass{stacks-project}
}{%
\documentclass{amsart}
}
\usepackage{amssymb}

% For dealing with references we use the comment environment
\usepackage{verbatim}
\usepackage{fancyvrb}
\usepackage{pmboxdraw}
\newenvironment{reference}{\comment}{\endcomment}
%\newenvironment{reference}{}{}
\newenvironment{slogan}{\comment}{\endcomment}
\newenvironment{history}{\comment}{\endcomment}

% For source figures retained as assets
\usepackage{graphicx}

% For commutative diagrams we use Xy-pic
\usepackage[all]{xy}

% We use 2cell for 2-commutative diagrams.
\xyoption{2cell}
\UseAllTwocells

% We use multicol for the list of chapters between chapters
\usepackage{multicol}

% This is generally recommended for better output
\usepackage{lmodern}
\usepackage[T1]{fontenc}
\usepackage{textalpha}

% For cross-file-references
% We prefer the snapshot tmp/xref/#2.aux,
% which latexmkrc replaces atomically, so
% that documents compiled simultaneously
% never read a half-written .aux file.
% The last argument keeps links to #2.pdf.
\usepackage{xr-hyper}
\newcommand{\maybeexternaldocument}[2]{%
  \IfFileExists{#2.aux}{%
    \externaldocument[#1]{#2}%
  }{%
    \PackageWarning{stack}{Missing #2.aux; external references with prefix #1 unavailable}%
  }%
}
\let\externaldocumentorig\externaldocument
\renewcommand{\externaldocument}[2][]{%
  \IfFileExists{tmp/xref/#2.aux}{%
    \externaldocumentorig[#1]{tmp/xref/#2}[#2.pdf]%
  }{%
  \IfFileExists{#2.aux}{%
    \externaldocumentorig[#1]{#2}%
  }{%
    \PackageWarning{stack}{Missing #2.aux; external references with prefix #1 unavailable}%
  }%
  }%
}

% Package for hypertext links:
\usepackage{hyperref}

% For any local file, say "hello.tex" you want to link to please
% use \externaldocument[hello-]{hello}
\externaldocument[introduction-]{introduction}
\externaldocument[conventions-]{conventions}
\externaldocument[sets-]{sets}
\externaldocument[categories-]{categories}
\externaldocument[topology-]{topology}
\externaldocument[sheaves-]{sheaves}
\externaldocument[sites-]{sites}
\externaldocument[stacks-]{stacks}
\externaldocument[fields-]{fields}
\externaldocument[algebra-]{algebra}
\externaldocument[brauer-]{brauer}
\externaldocument[homology-]{homology}
\externaldocument[derived-]{derived}
\externaldocument[simplicial-]{simplicial}
\externaldocument[more-algebra-]{more-algebra}
\externaldocument[smoothing-]{smoothing}
\externaldocument[modules-]{modules}
\externaldocument[sites-modules-]{sites-modules}
\externaldocument[injectives-]{injectives}
\externaldocument[cohomology-]{cohomology}
\externaldocument[sites-cohomology-]{sites-cohomology}
\externaldocument[dga-]{dga}
\externaldocument[dpa-]{dpa}
\externaldocument[sdga-]{sdga}
\externaldocument[hypercovering-]{hypercovering}
\externaldocument[schemes-]{schemes}
\externaldocument[constructions-]{constructions}
\externaldocument[properties-]{properties}
\externaldocument[morphisms-]{morphisms}
\externaldocument[coherent-]{coherent}
\externaldocument[divisors-]{divisors}
\externaldocument[limits-]{limits}
\externaldocument[varieties-]{varieties}
\externaldocument[topologies-]{topologies}
\externaldocument[descent-]{descent}
\externaldocument[perfect-]{perfect}
\externaldocument[more-morphisms-]{more-morphisms}
\externaldocument[flat-]{flat}
\externaldocument[groupoids-]{groupoids}
\externaldocument[more-groupoids-]{more-groupoids}
\externaldocument[etale-]{etale}
\externaldocument[chow-]{chow}
\externaldocument[intersection-]{intersection}
\externaldocument[pic-]{pic}
\externaldocument[weil-]{weil}
\externaldocument[adequate-]{adequate}
\externaldocument[dualizing-]{dualizing}
\externaldocument[duality-]{duality}
\externaldocument[discriminant-]{discriminant}
\externaldocument[derham-]{derham}
\externaldocument[local-cohomology-]{local-cohomology}
\externaldocument[algebraization-]{algebraization}
\externaldocument[curves-]{curves}
\externaldocument[resolve-]{resolve}
\externaldocument[models-]{models}
\externaldocument[functors-]{functors}
\externaldocument[equiv-]{equiv}
\externaldocument[pione-]{pione}
\externaldocument[etale-cohomology-]{etale-cohomology}
\externaldocument[proetale-]{proetale}
\externaldocument[relative-cycles-]{relative-cycles}
\externaldocument[more-etale-]{more-etale}
\externaldocument[trace-]{trace}
\externaldocument[crystalline-]{crystalline}
\externaldocument[spaces-]{spaces}
\externaldocument[spaces-properties-]{spaces-properties}
\externaldocument[spaces-morphisms-]{spaces-morphisms}
\externaldocument[decent-spaces-]{decent-spaces}
\externaldocument[spaces-cohomology-]{spaces-cohomology}
\externaldocument[spaces-limits-]{spaces-limits}
\externaldocument[spaces-divisors-]{spaces-divisors}
\externaldocument[spaces-over-fields-]{spaces-over-fields}
\externaldocument[spaces-topologies-]{spaces-topologies}
\externaldocument[spaces-descent-]{spaces-descent}
\externaldocument[spaces-perfect-]{spaces-perfect}
\externaldocument[spaces-more-morphisms-]{spaces-more-morphisms}
\externaldocument[spaces-flat-]{spaces-flat}
\externaldocument[spaces-groupoids-]{spaces-groupoids}
\externaldocument[spaces-more-groupoids-]{spaces-more-groupoids}
\externaldocument[bootstrap-]{bootstrap}
\externaldocument[spaces-pushouts-]{spaces-pushouts}
\externaldocument[spaces-chow-]{spaces-chow}
\externaldocument[groupoids-quotients-]{groupoids-quotients}
\externaldocument[spaces-more-cohomology-]{spaces-more-cohomology}
\externaldocument[spaces-simplicial-]{spaces-simplicial}
\externaldocument[spaces-duality-]{spaces-duality}
\externaldocument[formal-spaces-]{formal-spaces}
\externaldocument[restricted-]{restricted}
\externaldocument[spaces-resolve-]{spaces-resolve}
\externaldocument[formal-defos-]{formal-defos}
\externaldocument[defos-]{defos}
\externaldocument[cotangent-]{cotangent}
\externaldocument[examples-defos-]{examples-defos}
\externaldocument[algebraic-]{algebraic}
\externaldocument[examples-stacks-]{examples-stacks}
\externaldocument[stacks-sheaves-]{stacks-sheaves}
\externaldocument[criteria-]{criteria}
\externaldocument[artin-]{artin}
\externaldocument[quot-]{quot}
\externaldocument[stacks-properties-]{stacks-properties}
\externaldocument[stacks-morphisms-]{stacks-morphisms}
\externaldocument[stacks-limits-]{stacks-limits}
\externaldocument[stacks-cohomology-]{stacks-cohomology}
\externaldocument[stacks-perfect-]{stacks-perfect}
\externaldocument[stacks-introduction-]{stacks-introduction}
\externaldocument[stacks-more-morphisms-]{stacks-more-morphisms}
\externaldocument[stacks-geometry-]{stacks-geometry}
\externaldocument[moduli-]{moduli}
\externaldocument[moduli-curves-]{moduli-curves}
\externaldocument[examples-]{examples}
\externaldocument[exercises-]{exercises}
\externaldocument[guide-]{guide}
\externaldocument[desirables-]{desirables}
\externaldocument[coding-]{coding}
\externaldocument[obsolete-]{obsolete}
\externaldocument[fdl-]{fdl}
\externaldocument[index-]{index}

% Heap Project documents
\externaldocument[k3-]{k3}
\externaldocument[torelli-]{torelli}
\externaldocument[linear-algebra-]{linear-algebra}
\externaldocument[groups-]{groups}
\externaldocument[complex-analysis-]{complex-analysis}
\externaldocument[functional-analysis-]{functional-analysis}
\externaldocument[categories-]{categories}
\externaldocument[commutative-algebra-]{commutative-algebra}
\externaldocument[ringed-spaces-]{ringed-spaces}
\externaldocument[homological-algebra-]{homological-algebra}
\externaldocument[lie-algebras-]{lie-algebras}
\externaldocument[representation-theory-]{representation-theory}
\externaldocument[smooth-manifolds-]{smooth-manifolds}
\externaldocument[differential-topology-]{differential-topology}
\externaldocument[algebraic-topology-]{algebraic-topology}
\externaldocument[differential-geometry-]{differential-geometry}
\externaldocument[algebraic-geometry-]{algebraic-geometry}
\externaldocument[symplectic-geometry-]{symplectic-geometry}
\externaldocument[complex-geometry-]{complex-geometry}
\externaldocument[hodge-theory-]{hodge-theory}
\externaldocument[%
variations-of-hodge-structures-]{%
variations-of-hodge-structures}
\externaldocument[%
homological-mirror-symmetry-]{%
homological-mirror-symmetry}
\maybeexternaldocument{deformations-}{deformations}
\externaldocument[vertex-algebras-]{vertex-algebras}
\externaldocument[springer-]{springer}
\externaldocument[infty-categories-]{infty-categories}
\externaldocument[seminars-differential-geometry-]{%
seminars-differential-geometry}
\externaldocument[seminars-complex-geometry-]{%
seminars-complex-geometry}
\externaldocument[seminars-algebraic-geometry-]{%
seminars-algebraic-geometry}
\externaldocument[seminars-analysis-pde-]{%
seminars-analysis-pde}
\externaldocument[surfaces-]{surfaces}
\externaldocument[birational-maps-commutative-algebra-]{birational-maps-commutative-algebra}
\externaldocument[cremona-transformations-]{cremona-transformations}
\externaldocument[derived-categories-of-sheaves-]{derived-categories-of-sheaves}
\externaldocument[geometric-invariant-theory-]{geometric-invariant-theory}
\externaldocument[geometric-stability-conditions-and-group-actions-]{geometric-stability-conditions-and-group-actions}
\externaldocument[moduli-spaces-of-sheaves-]{moduli-spaces-of-sheaves}
\externaldocument[bridgeland-stability-general-theory-]{bridgeland-stability-general-theory}
\externaldocument[hilbert-schemes-]{hilbert-schemes}
\externaldocument[noncommutative-abelian-surfaces-and-kummer-type-hyperkahler-varieties-]{noncommutative-abelian-surfaces-and-kummer-type-hyperkahler-varieties}
\externaldocument[mumford-tate-groups-in-hodge-theory-]{mumford-tate-groups-in-hodge-theory}
\externaldocument[hodge-birational-atoms-2026-]{hodge-birational-atoms-2026}
\externaldocument[xxii-egd-2026-]{%
xxii-egd-2026}
\externaldocument[pde-]{pde}
\externaldocument[probability-]{probability}

\externaldocument[linguistics-]{linguistics}

\externaldocument[%
16th-alga-meeting-2026-]{%
16th-alga-meeting-2026}

\externaldocument[%
iv-jornadas-lefschetz-2026-]{%
iv-jornadas-lefschetz-2026}

\externaldocument[reading-the-masters-]{%
reading-the-masters}

\externaldocument[real-analysis-]{%
real-analysis}

\externaldocument[optimal-transport-]{%
optimal-transport}

\externaldocument[history-]{history}

\externaldocument[seiberg-witten-]{seiberg-witten}

\externaldocument[history-philosophy-biology-]{history-philosophy-biology}

\externaldocument[classical-mechanics-]{classical-mechanics}
\externaldocument[electromagnetism-]{electromagnetism}
\externaldocument[waves-and-optics-]{waves-and-optics}
\externaldocument[wif-]{waves-in-fluids-iv}
\externaldocument[thermal-physics-]{thermal-physics}
\externaldocument[quantum-mechanics-]{quantum-mechanics}
\externaldocument[special-relativity-]{special-relativity}
\externaldocument[physics-]{physics}

\externaldocument[chemistry-]{chemistry}

\externaldocument[molecular-cell-biology-]{molecular-cell-biology}

\externaldocument[microbiology-]{microbiology}

\externaldocument[protistology-]{protistology}

\externaldocument[mycology-]{mycology}

\externaldocument[botany-]{botany}

\externaldocument[zoology-]{zoology}

\externaldocument[animal-histology-]{animal-histology}

\externaldocument[animal-physiology-]{animal-physiology}

\externaldocument[systematics-biogeography-]{systematics-biogeography}

\externaldocument[ecology-natural-resources-]{ecology-natural-resources}

\externaldocument[earth-sciences-]{earth-sciences}

\externaldocument[biostatistics-]{biostatistics}

\externaldocument[mathematical-systems-biology-]{mathematical-systems-biology}

\externaldocument[cimpa-el-salvador-2024-]{cimpa-el-salvador-2024}

\externaldocument[anthropology-]{anthropology}

\externaldocument[yang-mills-theory-]{yang-mills-theory}

\externaldocument[seminars-algebra-]{seminars-algebra}

\externaldocument[seminars-physics-]{seminars-physics}

\externaldocument[seminars-computing-]{seminars-computing}

\externaldocument[oktobermat-2026-]{oktobermat-2026}

\externaldocument[saet-2026-]{saet-2026}

\externaldocument[seminars-symplectic-geometry-]{seminars-symplectic-geometry}

\externaldocument[seminars-dynamics-]{seminars-dynamics}

\externaldocument[oktobermat-2024-]{oktobermat-2024}

\externaldocument[quantum-groups-]{quantum-groups}

\externaldocument[legal-2027-]{legal-2027}

\externaldocument[derived-categories-]{derived-categories}

% Theorem environments.
%
\theoremstyle{plain}
\newtheorem{theorem}[subsection]{Theorem}
\newtheorem{proposition}[subsection]{Proposition}
\newtheorem{lemma}[subsection]{Lemma}
\newtheorem{corollary}[subsection]{Corollary}

\theoremstyle{definition}
\newtheorem{definition}[subsection]{Definition}
\newtheorem{example}[subsection]{Example}
\newtheorem{exercise}[subsection]{Exercise}
\newtheorem{situation}[subsection]{Situation}

\theoremstyle{remark}
\newtheorem{remark}[subsection]{Remark}
\newtheorem{remarks}[subsection]{Remarks}

\numberwithin{equation}{subsection}

% Macros
%
\def\lim{\mathop{\mathrm{lim}}\nolimits}
\def\colim{\mathop{\mathrm{colim}}\nolimits}
\def\Spec{\mathop{\mathrm{Spec}}}
\def\Hom{\mathop{\mathrm{Hom}}\nolimits}
\def\Ext{\mathop{\mathrm{Ext}}\nolimits}
\def\SheafHom{\mathop{\mathcal{H}\!\mathit{om}}\nolimits}
\def\SheafExt{\mathop{\mathcal{E}\!\mathit{xt}}\nolimits}
\def\Sch{\mathit{Sch}}
\def\Mor{\mathop{\mathrm{Mor}}\nolimits}
\def\Ob{\mathop{\mathrm{Ob}}\nolimits}
\def\Sh{\mathop{\mathit{Sh}}\nolimits}
\def\NL{\mathop{N\!L}\nolimits}
\def\CH{\mathop{\mathrm{CH}}\nolimits}
\def\proetale{{pro\text{-}\acute{e}tale}}
\def\etale{{\acute{e}tale}}
\def\QCoh{\mathit{QCoh}}
\def\Ker{\mathop{\mathrm{Ker}}}
\def\Im{\mathop{\mathrm{Im}}}
\def\Coker{\mathop{\mathrm{Coker}}}
\def\Coim{\mathop{\mathrm{Coim}}}

% Boxtimes
%
\DeclareMathSymbol{\boxtimes}{\mathbin}{AMSa}{"02}

%
% Macros for moduli stacks/spaces
%
\def\QCohstack{\mathcal{QC}\!\mathit{oh}}
\def\Cohstack{\mathcal{C}\!\mathit{oh}}
\def\Spacesstack{\mathcal{S}\!\mathit{paces}}
\def\Quotfunctor{\mathrm{Quot}}
\def\Hilbfunctor{\mathrm{Hilb}}
\def\Curvesstack{\mathcal{C}\!\mathit{urves}}
\def\Polarizedstack{\mathcal{P}\!\mathit{olarized}}
\def\Complexesstack{\mathcal{C}\!\mathit{omplexes}}
% \Pic is the operator that assigns to X its picard group, usage \Pic(X)
% \Picardstack_{X/B} denotes the Picard stack of X over B
% \Picardfunctor_{X/B} denotes the Picard functor of X over B
\def\Pic{\mathop{\mathrm{Pic}}\nolimits}
\def\Picardstack{\mathcal{P}\!\mathit{ic}}
\def\Picardfunctor{\mathrm{Pic}}
\def\Deformationcategory{\mathcal{D}\!\mathit{ef}}


\begin{document}

\title{Genealogía familiar}
\maketitle

\phantomsection
\label{section-phantom}

\tableofcontents

\section{Family history: evidence and scope}
\label{section-genealogia-mapa-general}

These notes concern Daniel's own family and
its connections with Mexican and Caribbean
intellectual history. Their public purpose is
documentary inquiry. Family recollections are
identified as such; neither a shared surname
nor an online tree establishes descent.

The following sections distinguish cited
biographical evidence, provisional
reconstructions, and recollections whose
speakers and dates still need recording. A
citation supports only the relationship it
actually describes. The reconstruction is
incomplete and is not a certified genealogy.

Three lines organize the inquiry: González
Casanova, Henríquez Ureña, and Lombardo
Toledano. The working reconstruction links Pablo
González Aznar and Encarnación Casanova
Castillo to Pablo González Casanova
(1889--1936), and Francisco Henríquez y
Carvajal and Salomé Ureña to Pedro Henríquez
Ureña. Pedro and Isabel Lombardo Toledano
lead to Natacha Henríquez Lombardo; Natacha
and Pablo González Casanova (1922--2023) form
the later family connection. These links must
be read with the evidence and outstanding
questions below, not as independently
verified edges of a tree.

Concepción del Valle Romo, the rector's
mother, was omitted from that abbreviated
diagram. Collateral branches and the author's
intervening generations are not reconstructed
here.

\section{Los González Casanova}
\label{section-genealogia-gonzalez-casanova}

Pablo González Casanova (1922--2023), rector de
la UNAM entre 1970 y 1972, fue hijo del lingüista
y filólogo Pablo González Casanova y de Concepción
del Valle Romo. Este parentesco aparece también en
las semblanzas institucionales de la UNAM.\footnote{%
UNAM, ``Valioso legado de Don Pablo'',
\url{https://www.cch.unam.mx/comunidad/valioso-legado-de-don-pablo}.}

En la generación anterior aparece la unión
de dos líneas: González y Casanova.
La reconstrucción provisional identifica
a Pablo González Aznar y Encarnación
Casanova Castillo como sus padres; falta
adjuntar el documento de este enlace.

La rama Casanova se reconstruye de modo
provisional, pendiente de
actas para cada enlace: José García-Casanova
y Rosa Vázquez; Francisco García-Casanova
Vázquez; Francisco García-Casanova Zaldívar;
Encarnación Casanova Castillo; Pablo González
Casanova (1889--1936); Pablo González
Casanova (1922--2023). La sucesión de nombres
conserva la hipótesis de trabajo anterior,
pero no demuestra cada relación de filiación.



La investigación pendiente importante es encontrar
el origen de José García-Casanova. El apellido
\emph{Casanova} existe tanto en Italia como en la
Península Ibérica; por eso no hay todavía base para
decir que esta rama sea italiana.
La reconstrucción anterior sitúa esta rama
en Yucatán en el siglo XIX; falta adjuntar
los documentos que permitan comprobarlo.

\section{Pedro Henríquez Ureña}
\label{section-genealogia-henriquez-urena}

Pedro Henríquez Ureña (1884--1946) nació en Santo
Domingo. Fue hijo de Francisco Henríquez y Carvajal
y de la poeta Salomé Ureña. Se casó con Isabel
Lombardo Toledano en México. Tuvieron dos hijas:
Natacha, nacida en México, y Sonia, nacida en
Argentina. Este último dato está confirmado por una
entrevista con Sonia Henríquez Ureña conservada por
El Colegio de México.\footnote{%
Mary Carmen Sánchez Ambriz, ``El acervo de Pedro
Henríquez Ureña a El Colegio de México. Entrevista
con Sonia Henríquez Ureña de Hlito'',
\url{https://libros.colmex.mx/wp-content/uploads/2021/04/boled_122.pdf}.}

La reconstrucción preliminar sitúa la
rama paterna de Pedro en Curazao y señala
a Noel Henríquez como abuelo. Falta citar
los documentos de cada enlace.
Hay una tradición familiar y
bibliográfica que relaciona esta rama con la
diáspora sefardí del Caribe y de los Países Bajos.
Éste es un punto que merece una investigación
archivística separada: antes de prolongar el árbol
hasta Europa quiero verificar los enlaces uno por
uno en registros de Curazao y, si es posible, en
archivos de comunidades sefardíes neerlandesas.

La línea materna de Pedro es la de Salomé Ureña,
una de las figuras fundamentales de la literatura y
la educación dominicanas. De este modo, esta rama
del árbol no es sólo genealógica: está directamente
ligada a la historia intelectual de la República
Dominicana y, por Pedro, a la historia intelectual
de México.

\section{Lombardo Toledano: la rama italiana}
\label{section-genealogia-lombardo-toledano}

Aquí sí aparece una procedencia italiana
explícitamente documentada. Vicente Lombardo
Toledano (1894--1968) fue hijo de Vicente Lombardo
Carpio e Isabel Toledano. Según la cronología
biográfica preparada por Rosa María Otero y Gama de
Lombardo, sus abuelos paternos fueron \emph{Vincenzo
Lombardo Cati}, originario de Turín, Italia, y
Marcelina Carpio. Vincenzo llegó a México en
1859.\footnote{%
Centro de Estudios Filosóficos, Políticos y Sociales
Vicente Lombardo Toledano, \emph{Escritos
autobiográficos}, cronología de Rosa María Otero y
Gama de Lombardo, p.~51,
\url{https://www.centrolombardo.edu.mx/%
wp-content/uploads/formidable/14/%
Escritos_autobiogr\%C3\%A1ficos_REV.pdf}.}

La misma fuente dice que los abuelos maternos de
Vicente eran Miguel Toledano y Flora Toledano,
descendientes de una antigua familia española cuyo
fundador había llegado a México en el siglo XVII.
Esto es especialmente interesante porque impide
identificar sin más ``Lombardo Toledano'' con una
sola procedencia: la rama Lombardo es italiana,
mientras que la rama Toledano está descrita allí
como española.

Isabel Lombardo Toledano, hermana de Vicente, se
casó con Pedro Henríquez Ureña.

La reconstrucción de esta rama enlaza a
Vincenzo Lombardo Cati y Marcelina Carpio con
Vicente Lombardo Carpio; a éste e Isabel
Toledano con Isabel Lombardo Toledano; y a
Isabel y Pedro Henríquez Ureña con Natacha.
La cronología citada documenta la ascendencia
de Vicente Lombardo Toledano. Los enlaces
adicionales de esta cadena requieren sus
propias referencias; no quedan probados por
esa sola cita.



El siguiente paso para esta rama es localizar el
acta de nacimiento de Vincenzo en el Piamonte y
reconstruir sus padres y abuelos. Sólo entonces se
podrá precisar más que ``originario de Turín''.

\section{Una red intelectual mexicana}
\label{section-genealogia-red-intelectual}

La genealogía se cruza de manera llamativa con la
historia intelectual mexicana. Pedro Henríquez Ureña
formó parte del núcleo del Ateneo de la Juventud
junto con figuras como Antonio Caso, Alfonso Reyes
y José Vasconcelos. Décadas después, Vicente
Lombardo Toledano fue alumno de Antonio Caso y se
convirtió en una figura central del debate político
y universitario mexicano.

Las relaciones entre Caso, Lombardo, Vasconcelos y
Henríquez Ureña fueron además personales e
institucionales, no sólo bibliográficas. La propia
cronología histórica de la UNAM registra, por
ejemplo, los conflictos universitarios de 1923
entre Lombardo y Vasconcelos, la renuncia de
Antonio Caso y la salida de Pedro Henríquez Ureña
de sus cargos docentes.\footnote{%
UNAM, cronología histórica, ``1920'', sección sobre
1923,
\url{https://www.unam.mx/acerca-de-la-unam/unam-en-el-tiempo/cronologia-historica-de-la-unam/1920}.}

Esto sugiere que una historia familiar completa no
debe ser sólo una lista de ascendientes. Hay también
una historia de redes: parentesco, amistad,
magisterio, instituciones y disputas intelectuales.

\section{Recuerdos familiares por verificar}
\label{%
section-genealogia-recuerdos-familiares}

Son recuerdos recogidos por Daniel.
Falta registrar informante, fecha y contexto
de cada uno. Se conservan como pistas de
investigación, no como hechos establecidos.

\begin{itemize}
\item En recuerdos familiares aparecen los nombres
Charles y Mara, Henrique y Lolita, Manuel y Helena,
y el apodo ``Titopaul'' para Pablo. Falta fijar con
precisión qué generación y qué parentesco designa
cada referencia.
\item Aparece también una ``tía Sonia''. Esto es
compatible con Sonia Henríquez Ureña, hermana de
Natacha, pero no debe identificarse automáticamente
sin contexto adicional.
\item Hay un recuerdo de una conexión familiar de
los Bracho--Carpizo--Valle con Antonio Caso. Todavía
no hay identificación documental suficiente para
dibujar esa rama.
\end{itemize}

\section{Preguntas abiertas}
\label{section-genealogia-preguntas-abiertas}

Las tres búsquedas principales que quedan abiertas
son:

\begin{enumerate}
\item \textbf{García-Casanova.} Encontrar el origen
de José García-Casanova y decidir documentalmente
si la línea anterior conduce a España, Italia u
otro lugar.
\item \textbf{Henríquez.} Reconstruir la línea de
Curazao generación por generación y comprobar la
tradición de ascendencia sefardí con registros
primarios.
\item \textbf{Lombardo.} Localizar a Vincenzo
Lombardo Cati en registros piamonteses y extender
la rama italiana hacia atrás.
\end{enumerate}

Estas notas deben actualizarse conforme aparezcan
actas civiles, parroquiales o documentos de archivo;
los árboles genealógicos en línea sirven como pistas,
no como prueba final.

\bibliographystyle{amsalpha}
\bibliography{my}

\end{document}
